\section{Kaggle lab: audit before modeling}
\label{sec:kaggle-titanic-audit}

\begin{kaggleproject}{Titanic --- data audit and simple feature construction}
\kagglesource{https://www.kaggle.com/datasets/yasserh/titanic-dataset}{Titanic Dataset}
\textbf{File.} \texttt{all\_datasets/titanic\_dataset/Titanic-Dataset.csv}\\
\textbf{Target.} \texttt{Survived}; binary classification later, but no model is fitted in this lab.\\
\textbf{Question.} What does one row represent, which columns are missing, how imbalanced is the target, and which transparent features can be constructed without using the target?

\textbf{Before running.} Predict which columns are likely to have missing values and which columns look like identifiers rather than model inputs.

\begin{labmode}
This is a highly guided foundations lab. Reproduce the reference audit faithfully first, then extend it with the requested checks and explanations.
\end{labmode}

\lstinputlisting[style=mlpython]{all_experiments/lab02-titanic_data_audit.py}

\textbf{What to inspect.} Do not hunt for a score. Inspect shape, dtypes, missingness, target prevalence, duplicated identities, and whether each engineered feature would exist at prediction time.

\begin{tryit}
Extend the audit without fitting a model. Add: (1) median age by passenger class, (2) survival rate by embarkation port, and (3) a duplicated-\texttt{PassengerId} check. For each calculation, write one sentence explaining what later decision it could influence.
\end{tryit}
\end{kaggleproject}


\begin{decisionmemo}
Define two prediction moments: beginning of term and late in term. Decide whether \texttt{G1} and \texttt{G2} are legitimate features for each moment. The correct feature set depends on the prediction-time information boundary, not on which version produces the highest score.
\end{decisionmemo}

\begin{labdeliverable}
Submit one audit notebook containing: a schema/missingness summary; the three requested Titanic checks; two prediction-time leakage decisions with justification; the Student Performance schema comparison; and a short note stating what you would need to know before merging or modeling the education files. No model score is required.
\end{labdeliverable}
